\Chapter{The Language of Predicate Logic}
\Lettrine[loversize=0.1,nindent=0pt]{Writing} clear, concise, and logical sentences is essential in mathematics. We need a consistent way to write statements to arrive at the axioms required for defining the foundations of mathematics. That is a job reserved for format language, a structure we will explore in this chapter for predicate logic.

\section{An Introduction to Formal Language}
In English, one may consider a sentence a sequence of words and think of an English word as a sequence of letters. However, not every sequence of letters is a valid word, and not every sequence of words is a grammatically sound sentence. The idea of formal language is to generalise this idea of valid words and valid sentences.

\begin{definition}[Formal Language]\index{formal language}
    A formal language $\Lang$ is a structure comprising two things.
    \begin{itemize}
        \item A set $\Alphabet$ of symbols, called the \term{alphabet} of $\Lang$.
        \item A \term{formal grammar} comprising \term{formation rules}\index{formation rule} that specify which \term{words}, which are sequences of symbols from $\Alphabet$, belong to $\Lang$.
    \end{itemize}
\end{definition}

We will explore more about alphabets, words, formal grammar, and their related constructs in the following pages.

\subsection{Alphabets and Strings}
In English, our (primary) alphabet is A to Z, along with the two cases each letter can come in (i.e., uppercase or lowercase). These 54 characters define all the valid letters we can use in English. We use these valid letters to form words via concatenating letters together. Likewise, we can define symbols that make up our alphabet in a formal language and define strings that use the symbols present in our alphabet.

\begin{definition}[Alphabet]
    Let $\Lang$ be a formal language. The \term{alphabet}\index{alphabet}\index{formal language!alphabet} of $\Lang$, usually denoted $\Alphabet$, is the set of symbols used to construct strings in $\Lang$.
\end{definition}

\begin{definition}[String]
    Let $\Lang$ be a formal language. A \term{string}\index{string}\index{formal language!string} (or string) is either a finite concatenation of symbols from $\Alphabet$ in $\Lang$ or the empty string $\emptystring$.

    The \term{length}\index{string!length} of a string is the number of symbols making up the string, with the empty string having length 0.
\end{definition}
\begin{remark}[Making Strings]
    Some authors may permit strings to be formed from other means, not just concatenation. For example, using $\Alphabet = \{a, b\}$, a string could be formed from other string by the rule
    \[
        \square b,
    \]
    i.e. all words must end with $b$, or
    \[
        aaa\square bab \square aaa,
    \]
    i.e. two (not necessarily distinct) strings can be joined together with $bab$ in the middle and by placing $aaa$ at the front or end of the combined string. However we will stick with simple concatenation.
\end{remark}

\begin{example}[A Simple Alphabet]
    Let $\Alphabet = \{a, b\}$. Then $\emptystring$ (the empty string), $a$, $b$, $aa$, $ab$ $ba$, $bb$, $aaa$, $aab$, $aba$, $baa$, etc. are all valid strings. Notice that we have not defined any rules on what strings belong in our formal language -- these are just strings formed from $\Alphabet$.
\end{example}

\subsection{Grammar and Well-Formed Strings}
TODO: Add
% Add remark that 'words' here means 'well-formed strings'

\section{Language of Predicate Logic}
TODO: Add

\subsection{Variables and Formulae}
TODO: Add

\subsection{Grammar and Well-Formed Formulae}
TODO: Add
